% SOURCE docs/release/2026-10-29/source/docs/writing-repair/2026-10-29/src/aeceff940e6ca41b.tex CHAR 0
\chapter{Variables, diagrams and justified disagreement}
\label{la:diagrams}

The preceding chapter established what definitions and fallible evidence can
contribute. Here those objects enter a reviewable diagram. The argument first
separates graph classes, premises and retrieval outcomes, then develops exact
proxy, collider, clock and disagreement examples. The ontology comparison
connects the diagram to implementable treatment versions and source records.
The final restrictions study follows a premise through both graph and effect
sets, showing why a smaller graph class need not settle the decision.

\section{From vocabulary to measured variables}

A useful causal diagram begins with variables whose meanings can survive
contact with data. A business glossary may contain ``risk,'' ``relationship''
and ``engagement,'' while a dataset contains a risk-score vintage, a set of
linked accounts and an observed click. Those are related concepts, but they
are not interchangeable causal variables. An assistant can propose mappings,
detect inconsistent names and assemble definitions. Each proposed mapping
still needs an operational meaning, a unit and an information clock.

OntoLearner offers relevant modular ideas for ontology loading, typing,
taxonomy and relation-learning tasks. The inspected paper reports
domain-dependent errors and retrieval failures. Its relation-learning
capability is not a guarantee that a relation is causal, or that a concept
has a valid measured proxy in the bank. Ontology alignment should therefore
precede and inform graph construction without being silently promoted into
causal identification. \cite{v12_la_2607_01977}

For example, ``customer owns account'' may describe a legal relation,
``records share an identifier'' an administrative relation, and ``records
have a high linkage score'' a model output. If ownership affects both offer
eligibility and repayment behavior, substituting the linkage score for
ownership leaves a measurement problem. Adding an arrow from the score to
repayment does not resolve it. The diagram should distinguish latent
ownership, observed records, the linkage procedure and sample selection.

Likewise, an ordinal description such as ``high influence'' cannot be used
as an effect size without a scale. An assistant's rating of a proposed edge
may express a preference among textual explanations. It is not automatically
a probability that the edge is present, a posterior over graphs, or a
confidence level. A typed diagram can retain such a rating as explicitly
uncalibrated metadata while refusing to use it as scientific evidence.

\section{What an edge asserts}

An edge proposal should state more than two variable names. The endpoint
marks identify the proposed graphical object. A source and inspected anchor
state what supports the proposal. A scope describes the population, period
and intervention to which that evidence applies. A counterevidence record
states whether contrary evidence was found, absent in the inspected material,
or not successfully searched. A review disposition records acceptance as a
premise, rejection, or unresolved status.

These fields serve different purposes. A source can be authentic but
inapplicable to the current banking action. A study can concern a different
population or a different definition of treatment. A plausible timing argument
can constrain direction without establishing that any causal connection
exists. The assistant should make the supported proposition as narrow as
the evidence warrants and retain the reasons for disagreement.

CausalChat supplies an example of interactive model development: the inspected
paper combines optional CPDAG discovery, model suggestions, graph editing
and structural-equation coefficient updates. Its prompts use ordinal
causal-significance ratings and ask for possible confounders and mediators.
These are useful interaction ideas. The paper's expert feedback also raises
missing citations and uncertainty about the quality of scientific evidence
provided by the model. A coefficient update after adding an edge cannot
establish that the edge itself is true. \cite{v12_la_2410_14146}

The proposed banking workflow retains the interaction but changes the
scientific handoff. An accepted premise is labeled as a premise with a
responsible review, not as a truth certified by the model. A model proposal
alone cannot remove an observationally compatible graph alternative. The
tool responsible for formal identification receives the complete reviewed
graph object and its restrictions, rather than a picture whose uncertainty
has disappeared in editing.

\section{Graph classes and collections of candidate models}

A CPDAG or PAG is a structured representation with its own interpretation.
A list of proposed DAGs is another object. A collection of edges marked
``possible'' is not automatically an enumeration of every combination of
those edges. Confusing these representations can silently exclude models
or include incompatible ones.

The project contract preserves a CPDAG or PAG directly when that is the
appropriate object, including circle endpoints. It also permits explicitly
listed candidate models. An edge marked possible within a candidate describes
the evidence status of that candidate's edge; it does not make all optional
edge combinations an exhaustive graph class. Downstream calculations must
know which interpretation applies.

An exact binary example shows the consequence of retaining alternatives.
In model \(F\), a binary offer \(A\) is Bernoulli with probability one half,
and response \(Y\) equals \(A\) with probability \(0.9\). In model \(R\),
\(Y\) is Bernoulli with probability one half, and \(A\) equals \(Y\) with
probability \(0.9\). Both models produce
\[
\begin{array}{c|cc}
 &Y=0&Y=1\\\hline
A=0&9/20&1/20\\
A=1&1/20&9/20
\end{array}.
\]
In \(F\), setting \(A\) from zero to one changes \(E[Y]\) by \(0.8\).
In \(R\), intervening on the equation for \(A\) leaves \(Y\) unchanged.
The two compatible constructions therefore imply the effect candidates
\(\{0,0.8\}\). A directional suggestion favoring \(F\) does not provide
new identifying information unless its underlying premise is independently
justified. These finite calculations are project constructions, not a
calibrated confidence set for a learned graph.

More generally, for a declared graph class \(\mathcal G\), one relevant
downstream object is
\begin{equation}
 \Theta(P,Q,\mathcal G)=\bigcup_{G\in\mathcal G}\Theta_G(P,Q),
 \label{la:graph-union}
\end{equation}
where \(\Theta_G(P,Q)\) contains effects compatible with distribution \(P\),
question \(Q\) and the model restrictions for \(G\). A nonidentified member
can make the union broad. A list of estimated points is not automatically
the whole union, and graph-selection uncertainty cannot be replaced with
ordinary coefficient uncertainty after selecting one member.

Equation~\eqref{la:graph-union} also clarifies what a disagreement review
should resolve. If two experts differ only on an edge that leaves the target
functional unchanged, that disagreement may be immaterial for the current
decision while still mattering for another target. If their graphs imply
different effect signs or feasibility, the conflict is decision-relevant.
The assistant can identify that consequence and prioritize the required
evidence. It should not average expert preferences into an unexplained graph.

\section{Retrieval failure and scientific disagreement}

An evidence assistant can fail to find a passage because the document is
unavailable, parsing failed, the query was poor, the retriever missed it, or
the passage is absent from the inspected document. These cases support
different statements. ``No passage was retrieved'' cannot by itself become
``the authors ignored the assumption,'' and neither statement establishes
that the assumption is false in the population.

The inspected ARGUS paper makes this distinction especially useful. It
audits reported evidence concerning DiD identifying assumptions and does
not claim to adjudicate causal truth. Its synthetic flaw study exhibits a
retrieval/abstention tradeoff, and its small human-label pilot reports an
in-sample calibration change that was not tested on unseen papers. These
observations motivate a banking evidence state that is separate from a
scientific risk judgment. \cite{v12_la_2609_30867}

The proposed record contains the source bytes, the query or retrieval
configuration, inspected anchors and the resulting narrow claim. If evidence
is missing, it records the missing step. If two sources conflict, it records
their populations, dates and definitions before treating the conflict as
substantive. The analyst can then determine whether the sources concern
different regimes or whether a genuine contradiction remains.

This process does not require pretending that every edge has a published
study behind it. Some causal premises depend on institutional facts: which
system sets a rate, when a policy becomes visible, how an audit was selected,
or which identifier establishes ownership. Those facts need an accountable
owner and authentic records. Literature retrieval can help formulate the
questions, but it cannot supply an institution's missing assignment truth.

\section{A reviewable editing sequence}

The proposed interaction starts with a question and variable definitions,
then records diagram alternatives. Each edit changes a declared object:
add or remove a variable, refine a clock, replace a proxy, propose an endpoint,
or accept a restriction. The edit retains its predecessor identity and
reason. A reviewer can therefore see whether an apparent improvement in
fit came from a new assumption, a different target, or an actual correction
of an error.

Formal checks can reject a directed cycle in a DAG, impossible declared
time order or an explicitly forbidden direction. They can flag mediator,
collider or common-ancestor patterns in the supported representation. Such
checks help maintain a coherent declaration but do not cover every graph
class or establish that an accepted edge is causally true. In particular,
a diagram with no syntax error can still omit an important common cause or
misrepresent the observation process.

The user should be able to retain an unresolved diagram and ask for analyses
conditional on several alternatives. A refusal to pick one graph can be the
scientifically correct output. The later discovery and identification
operations then need to report what changes across those alternatives and
what remains common. This makes uncertainty operational: it directs the
next information request rather than appearing only as a disclaimer after
a single confident answer.

The following constructions make that consequence explicit. They separate
observational equivalence from sampling precision, measured proxies from
confounders, and a graph edit from evidence that warrants it. Their monetary
comparisons then locate the disagreements worth investigating.


% SOURCE docs/release/2026-10-29/source/docs/writing-repair/2026-10-29/src/d152488cfd496484.tex CHAR 0
\section{Equal observed laws can imply different intervention answers}
\label{v13:la-equivalence}

A diagram assistant must recognize when available observations cannot determine
the requested direction or effect. Consider two project-constructed binary
models. Let $U$ be fair Bernoulli and $A=U$. In model one $Y=A$; in model two
$Y=U$. Both observed datasets contain only $(A,Y)=(0,0)$ and $(1,1)$ with equal
probability. Under an intervention setting $A=a$, model one has mean outcome
$a$ while model two has mean outcome $1/2$. Their average effects are one and
zero. Exact enumeration establishes the discrepancy. \cite{v13_project_exact}

Estimating the observed marginal propensity as $1/2$ does not solve the
problem: exchangeability given the available information is precisely what is
missing. A model that recognizes a plausible causal story may prefer one graph,
but that preference is not identification from these data. Expert knowledge,
measurements or experiments could distinguish the worlds; those inputs need
their own evidence and should remain visible in the result.

A continuous version shows why effect sets can be unbounded. Let independent
standard Gaussian variables $U$ and $\varepsilon$ generate
\[
 X=U,\qquad Y=bX-bU+\varepsilon.
\]
The factual law is $X=U,Y=\varepsilon$ for every real $b$. Under
$\operatorname{do}(X=x)$, however, the mean of $Y$ is $bx$. Therefore every
real intervention slope is compatible with the same observed distribution
within this allowed model class. More observations estimate that common law
more accurately; they do not select $b$. A discovery tool or language model
cannot remove this ambiguity without an additional restriction.

This is a useful evaluation case for agentic discovery. A high-quality response
retains the competing models, explains their observational equivalence and
identifies the information that could separate them. An overconfident single
graph should be scored as a substantive error even if its prose is convincing.
A refusal or whole-line effect set can be the scientifically correct output;
its lack of immediate decision usefulness should be reported separately.

\section{Prior knowledge needs a typed evidence record}

The inspected DirectLiNGAM interface distinguishes a directed path, no directed
path and unspecified prior knowledge. The y0 identification dispatcher requires
a supplied graph rather than eliciting its truth. These interfaces make the
translation from narrative evidence to formal restrictions consequential.
An agent saying ``A is likely relevant to Y'' has not supplied a justified
directed-path constraint. \cite{v13_sw_lingam,v13_sw_y0}

For each proposed restriction, a useful record distinguishes a temporal order,
an intervention mechanism, a measured association, a domain assertion and a
formal exclusion. It binds the variables, their measurement/availability clocks,
the source passage, the supporting scope and any disagreement. A source about
one population or intervention need not justify a hard exclusion in another.
The record should preserve which conclusions change if the assertion is removed.

Suppose a manager states that a rate change cannot be caused by next month's
withdrawal. That supports a temporal ordering when the rate decision was fixed
before the withdrawal. It does not rule out an earlier expectation of withdrawal
affecting the rate, or a common market shock affecting both. A diagram assistant
can expose the omitted expectation variable and the common-cause alternative.
It should not convert the temporal statement into unconfoundedness.

An analogous issue arises in network diagrams. MossSpider's selected source
uses a directed influence convention that requires explicit checking, while a
bank may draw arrows from guarantor to beneficiary. A language model can produce
the correct-looking graph but an incorrectly oriented adjacency representation.
A two-node semantic test with known distinct attributes is more informative
than a visual similarity score. \cite{v13_sw_mossspider}

\section{Set-valued reasoning should survive verbal reporting}

If uncertainty concerns several graphs or nuisance parameters, the effect
object may be a union of intervals. A project ratio example illustrates the
boundary. Let a numerator lie in $[1,2]$ and an audit multiplier in
$[-1/2,1/2]$, with the relation $m=r\theta$. Excluding division by zero while
retaining all feasible nonzero multipliers gives
\[
 \Theta=(-\infty,-2]\ \cup\ [2,\infty).
\]
The gap around zero and the unbounded tails are substantive. Replacing this
object by its hull loses the exclusion of zero; replacing it by a finite
interval loses compatible effects. A monetary transformation should be applied
to each component and preserve endpoint inclusion and infinite endpoints.

The agent should describe why the set is disconnected, what assumptions support
it and what information would make it more useful. It should not translate
``wide uncertainty'' into an invented central estimate. When both zero numerator
and zero multiplier remain jointly feasible, the effect can instead be the
whole real line. The distinction depends on the joint region, not just
separate point estimates.

This reasoning also separates model plausibility from sampling coverage. If a
joint confidence region contains the true observable and nuisance parameters,
its effect image contains the true effect on that event. The implication is
elementary; constructing the valid joint region under the bank's actual
dependence and audit law is the difficult statistical work. An agent can
propagate a supplied region faithfully. It cannot infer that its coverage is
valid merely because the downstream set arithmetic passed a test.


% SOURCE docs/release/2026-10-29/source/docs/writing-repair/2026-10-29/src/4741d08031d9bab0.tex CHAR 0
\hypertarget{comparing-four-ways-to-assist-diagram-construction}{%
\section{Comparing four ways to assist diagram
construction}\label{comparing-four-ways-to-assist-diagram-construction}}

The useful comparison is between the operations an assistant performs
and the evidence each operation adds. An interactive editor helps an
analyst see and revise a hypothesis. A language model can propose
mechanisms and missing variables from descriptions. A response-based
conditional-independence procedure can organize model judgments into a
graph search. A workflow agent can select and operate a statistical
discovery method under declared assumptions. These operations may share
an interface, but they answer different scientific questions. The
proposed product should preserve that distinction when it combines them.

CausalChat illustrates the interactive route. Its inspected preprint
describes an optional CPDAG obtained from data, GPT-4 comments on
relations, visual summaries linked to model responses, and
structural-equation coefficient updates after graph changes. This is a
useful design for making an analyst's reasoning inspectable. Its
evidence for interface benefits nevertheless needs careful
interpretation. Six university students first resolved three
self-selected edges using BIC feedback, then another three using LLM
text, and finally the remaining three using CausalChat Lite. The
reported error counts were 5/18, 7/18 and 1/18. The edge sets and order
changed together with the interface. Those results describe the authors'
experiment; they do not identify the effect of replacing one interface
with another on matched, randomly assigned banking tasks. Expert
participants also questioned missing citations and the quality of the
scientific evidence behind model suggestions. A compact visualization
may improve comprehension while leaving the causal premise unresolved.
\citep{v12_la_2410_14146}

The breadth-first method of Jiralerspong and colleagues offers a
different contribution. A query expands one variable into its proposed
effects, and insertion rejects an edge that would create a directed
cycle. The authors' linear query-count argument follows from visiting
each variable once. It does not make the total input/output token count
linear: each expansion can refer to many variables and the accumulated
graph. The reported 221-node experiment achieved an F1 score of 0.351,
while some comparison configurations and the proposed variant using
observational statistics were omitted for resource reasons. That
evidence motivates an efficiency baseline with actual tokens, latency
and omitted requests accounted for. It supplies no reason to certify
every retained arrow merely because the final graph is acyclic.
\citep{v12_la_2402_01207}

chatPC places language-model judgments inside a conditional-independence
search. Its prompts include variable descriptions, context and a
conditioning set; its repeated-answer test concerns the relative
probabilities of YES and NO responses. The paper reports agreement under
exchange of the two queried variables of 0.510 for GPT-3.5 and 0.870 for
GPT-4, with the latter often rejecting independence. Symmetry is a
necessary property of a conditional-independence relation. Measuring the
frequency with which a model respects it is useful, but that frequency
is not the error rate of a statistical test in the bank's observation
law. The two sources of randomness must remain separate: a model can
consistently repeat a mistaken judgment about the world.
\citep{v12_la_2406_07378}

The causal-learn+ system description instead places agents around formal
discovery: data inspection, assumption elicitation, method
recommendation, tool coordination and interpretation. It separates
suggestions, user decisions, algorithm settings and graph outputs, and
emphasizes retaining CPDAG and PAG ambiguity. This is a strong
architectural alternative when the aim is access to an existing formal
method. Its call for future benchmarks of diagnosis, selection and
calibrated interpretation also limits the performance conclusions
available from the description itself. Keeping a formal algorithm at the
center does not establish its premises, and user approval does not turn
a disputed premise into empirical truth. Conversely, a fully specified
and evaluated hybrid method is conceivable; the architecture's
restrictive preference is not an impossibility theorem about every use
of language-model information. \citep{v12_la_2606_23608}

These comparisons support a product with several explicitly named
operations. Proposing a missing common cause should return a hypothesis
and its evidence. Editing a reviewed graph should return a new graph and
the changed premises. Running a specified discovery method should retain
the method's actual output class. Identifying a target should consume
that class through a supported formal path. A user can benefit from all
four services without being led to believe that a plausible graph
proposal has already passed the last two.

\hypertarget{a-measured-risk-score-is-not-the-risk-variable-it-names}{%
\section{A measured risk score is not the risk variable it
names}\label{a-measured-risk-score-is-not-the-risk-variable-it-names}}

Variable elicitation is consequential before the first edge is drawn.
Consider a stipulated servicing example. Let \(U\) be a customer's
unobserved financial resilience, \(A\) an assistance offer, \(Y\)
successful repayment, and \(Z\) a pre-offer score category. Take all
variables to be binary. The generated mechanisms are

\[
 P(U=1)=\tfrac12,\qquad
 P(A=1\mid U=u)=\tfrac15+\tfrac35u,
\]

\[
 P(Y=1\mid A=a,U=u)=\tfrac1{10}+\tfrac15a+\tfrac12u,
 \qquad P(Z=u\mid U=u)=\tfrac45.
\]

The score error is independent of the action and outcome mechanisms
conditional on \(U\). The offer is more frequent in the more resilient
category in this stipulated assignment law. The graph contains
\(U\to A\), \(U\to Y\), \(U\to Z\) and \(A\to Y\). A dictionary that
calls both \(U\) and \(Z\) ``customer risk'' can conceal the measurement
distinction without changing a single recorded number.

Under an intervention on \(A\), the distribution of \(U\) remains one
half in each category. The mean success probabilities are

\[
 E\{Y(0)\}=\tfrac7{20},\qquad
 E\{Y(1)\}=\tfrac{11}{20},\qquad
 \tau=\tfrac15.
\]

In the observed groups, however, \(P(U=1\mid A=1)=4/5\) and
\(P(U=1\mid A=0)=1/5\). Their observed difference is therefore

\[
 E(Y\mid A=1)-E(Y\mid A=0)
 =\tfrac15+\tfrac12\left(\tfrac45-\tfrac15\right)=\tfrac12.
\]

Adjustment for the noisy score improves neither the scientific premise
by declaration nor the graph by relabeling. For \(Z=1\), Bayes' rule
gives

\[
 P(U=1\mid A=1,Z=1)=\tfrac{16}{17},\qquad
 P(U=1\mid A=0,Z=1)=\tfrac12.
\]

For \(Z=0\), the corresponding probabilities are \(1/2\) and \(1/17\).
The residual difference is \(15/34\) in either score stratum, so the
standardized contrast is

\[
 \sum_z\{E(Y\mid A=1,z)-E(Y\mid A=0,z)\}P(z)
 =\tfrac15+\tfrac12\tfrac{15}{34}
 =\tfrac{143}{340}\simeq .421.
\]

It remains above the true effect by \(15/68\), despite using a variable
measured before the offer. All sixteen cells of the joint generated law
of \((U,A,Z,Y)\) have positive probability. This is not a failure caused
merely by an empty observed stratum. It is the consequence of treating a
proxy as if it blocked the latent common-cause path. Exact enumeration
reproduces each calculation. \citep{exec13_diagram_examples}

An assistant can help by asking what the score measures, who constructed
it, what data were available at scoring time and whether the assignment
policy used the score or additional information. It can also propose
alternative diagrams: the offer may depend only on the recorded score,
or it may depend on latent information also seen by a servicing officer.
Those alternatives have different adjustment implications. A
language-model statement that the score is ``a confounder'' is too
coarse to settle the question.

For a proposal, publicly available dictionaries and controlled examples
are sufficient to demonstrate this failure and specify the required
interface. They cannot establish the target bank's true measurement
process. The immediate design should retain the conceptual variable,
measured proxy and measurement mechanism separately. An eventual bank
study of the proxy is justified only if the prototype and decision
analysis show that the unresolved measurement information is
consequential.

\hypertarget{what-additional-measurement-information-would-buy}{%
\section{What additional measurement information would
buy}\label{what-additional-measurement-information-would-buy}}

The proxy example also gives a concrete route beyond a generic request
for better data. Suppose a symmetric error probability \(e\) is known,
and the stronger nondifferential measurement condition

\[
 P(Z\mid U,A,Y)=P(Z\mid U)
\]

holds. For fixed \((a,y)\), write \(o_z=P(A=a,Y=y,Z=z)\) and
\(q_u=P(A=a,Y=y,U=u)\). Then

\[
 \begin{pmatrix}o_0\\o_1\end{pmatrix}
 =
 \begin{pmatrix}1-e&e\\e&1-e\end{pmatrix}
 \begin{pmatrix}q_0\\q_1\end{pmatrix}.
\]

The determinant is \(1-2e\). For the declared range \(0\leq e<1/2\),
inversion yields

\[
 q_u=\frac{(1-e)o_u-eo_{1-u}}{1-2e}.
\]

Thus known measurement information and the stated mechanism permit
reconstruction of the latent joint law in this finite example.
Adjustment using that reconstructed \(U\) law recovers \(1/5\), exactly
matching the intervention calculation. No observed customer identity or
latent label has been manufactured: the reconstruction depends
explicitly on an assumed, known error matrix.
\citep{exec13_diagram_examples}

The route has a meaningful failure boundary. In the vector one-norm, the
inverse matrix has norm \(1/(1-2e)\), because each column's absolute
entries sum to that value. At \(e=1/5\), a worst-case perturbation in
the observed joint probabilities can be amplified by \(5/3\) before the
further ratios needed for adjustment. As \(e\) approaches one half,
measurement inversion becomes unstable. At one half, the score supplies
no information about which latent category generated it, and this
inversion is unavailable. An exact symbolic inverse alone does not make
the resulting effect precise.

In finite samples the inverted entries can become negative. Such entries
may reflect sampling error, a misspecified measurement model or an error
matrix incompatible with the observed law. Clipping them to zero and
renormalizing changes the estimator and its interpretation. A defensible
alternative is to retain a constrained uncertainty region for the
observed probabilities and measurement parameters, then project the
compatible latent laws into effect values. Whether that region has
justified coverage is an additional statistical obligation; the exact
population example does not supply it.

If \(e\) is estimated from a labeled subset, the subset's selection and
label truth matter. A convenience sample of easily resolved customers
need not identify the same error matrix for difficult cases. If
difficulty affects the error rate even after conditioning on \(U\), the
symmetric matrix is wrong. If the score was constructed using post-offer
outcomes, nondifferential measurement can fail. The assistant's useful
contribution is to identify these propositions, locate relevant
documentation and propose a sensitivity analysis. Repeating a confidence
rating about the score cannot replace the missing information.

This makes the next experiment more specific. A public or generated
benchmark can compare inversion, naive proxy adjustment and a valid
refusal as the measurement error changes. It can establish that the
proposed software preserves the declared assumptions and numerical
limits. A later target-bank measurement audit would answer a different
question: whether an adequate and sufficiently representative
measurement model exists there. The expense of that audit belongs after
evidence that the affected decision warrants it.

\hypertarget{post-action-engagement-can-be-a-collider}{%
\section{Post-action engagement can be a
collider}\label{post-action-engagement-can-be-a-collider}}

The label ``engagement'' often invites adjustment, because it sounds
like a useful customer characteristic. Its clock and causal role matter
more than its name. Consider another exact example with independent fair
binary variables \(A\) and \(U\). Let \(A\) be a randomized outreach
offer, \(U\) latent willingness to repay, \(Y=U\), and \(M\) a
post-offer engagement indicator satisfying

\[
 P(M=1\mid A=a,U=u)=\tfrac1{10}+\tfrac25a+\tfrac25u.
\]

The graph has \(A\to M\leftarrow U\to Y\). The offer has zero causal
effect on repayment, and the unadjusted observed difference is zero. All
eight observed cells of \((A,M,Y)\) have positive probability.
Conditioning on \(M\) nevertheless makes the randomized action
informative about \(U\).

Among engaged customers,

\[
 P(Y=1\mid A=1,M=1)=\frac{.9}{.5+.9}=\tfrac9{14},
 \qquad
 P(Y=1\mid A=0,M=1)=\frac{.5}{.1+.5}=\tfrac56.
\]

Their difference is \(-4/21\). Among customers with \(M=0\), the
corresponding probabilities are \(1/6\) and \(5/14\), giving the same
difference. Standardizing over \(M\) therefore reports
\(-4/21\simeq-.190\) even though the true effect is zero. The apparently
richer adjustment set has created the discrepancy.
\citep{exec13_diagram_examples}

The example separates three questions that an assistant should not
merge. The total effect concerns \(E\{Y(1)-Y(0)\}\). An observed
comparison among engaged customers conditions on \(M=1\) after the
action and can compare different latent populations across arms. A
principal-stratum question instead refers to a joint potential-response
condition, such as customers who would engage under either offer
assignment. Its definition and identifying assumptions are different
again. A clearer English description does not make those targets
interchangeable.

If engagement is a mediator in another plausible graph, holding it fixed
can remove part of the pathway through which the offer works. A
controlled direct effect would involve \(Y(a,m)\) and an intervention on
engagement, if that intervention has a coherent meaning. Ordinary
regression adjustment for an observed mediator does not automatically
identify that effect, especially with hidden causes of both the mediator
and repayment. The assistant should first expose the target choice and
then ask which scientific conditions support the requested calculation.

The remedy is not a universal ban on using any variable recorded after
the offer. A genuinely baseline quantity may be entered into a database
later, and retrospective measurement can be relevant to an analysis if
its measurement and selection assumptions are justified. Conversely, a
variable entered early may encode a forecast produced using information
unavailable to the policy at the decision time. The product needs
distinct event, measurement and availability clocks, together with an
account of how the variable was constructed.

\hypertarget{time-unrolling-and-the-difference-between-analysis-and-deployment}{%
\section{Time unrolling and the difference between analysis and
deployment}\label{time-unrolling-and-the-difference-between-analysis-and-deployment}}

An institution may describe a feedback loop: offers affect balances, and
balances affect later offers. A graph over two timeless variables,
``offer'' and ``balance,'' can then contain a directed cycle. That
observation does not establish that every causal model of the workflow
requires a cyclic graph. With dates made explicit, a possible structure
is

\[
 A_t\longrightarrow B_{t+1}\longrightarrow A_{t+1}
 \longrightarrow B_{t+2}.
\]

The arrows describe distinct variable instances. An acyclic
time-unrolled model can represent feedback across decisions while
retaining a clear within-date ordering. Whether it is scientifically
adequate still depends on the sampling frequency, contemporaneous
mechanisms, hidden information and the policy that chooses each action.
Deleting an inconvenient arrow from the timeless graph would conceal
these issues rather than resolve them.

For each instance, the proposed variable record distinguishes the period
described, the event that generated its value, the time it was first
available and any subsequent correction. A customer's end-of-month
balance may be known only after settlement. A revised macroeconomic
release can describe an earlier month while containing information that
was not available when an offer was chosen. Using the revised value to
reconstruct a historical policy can create information leakage without
introducing an obviously backward arrow into a diagram.

Write \(\mathcal H_t\) for the information actually available when an
action is selected. A deployable policy has the form
\(A_t=\pi_t(\mathcal H_t)\), perhaps with explicit randomized variation.
A retrospective adjustment variable need not belong to exactly the same
information set: it might help describe pre-action confounding after
being measured later. The scientific validity of that retrospective use
and the feasibility of using it in a prospective policy are separate
questions. The diagram assistant should preserve both records rather
than infer feasibility from a successful estimation call.

A bank also needs to distinguish a true event from the observation of
that event. An offer can affect whether a customer responds to a survey,
whether a missing identifier is resolved or whether an account enters
the analysis extract. If \(S\) denotes observation, an effect among
\(S=1\) can concern a selected population whose composition changes with
the action. Adding an ``observed'' flag to a data table is not enough;
its parents belong in the causal model. Public synthetic benchmarks can
vary this mechanism deliberately and reveal whether the interface
silently drops the selection node.

These distinctions guide the assistant's questions. A temporal
contradiction should generate a request to refine the variable instance
or clock. A change in data availability should affect the deployable
policy record. An unresolved observation mechanism should remain a
possible source of selection bias. None requires the assistant to invent
a hidden construct or declare a fully specified longitudinal
identification result before the required evidence exists.

\hypertarget{evidence-for-an-arrow-a-path-and-an-exclusion}{%
\section{Evidence for an arrow, a path and an
exclusion}\label{evidence-for-an-arrow-a-path-and-an-exclusion}}

Expert statements differ in logical strength. ``The offer precedes the
withdrawal'' can rule out a particular backward-in-time mechanism under
the declared clock. It does not establish an arrow from offer to
withdrawal, exclude an earlier common cause or prove a nonzero effect.
``The offer changes the fee, which changes withdrawal'' proposes a
directed path through a mediator; it does not necessarily propose a
direct arrow after that mediator is included. ``The offer was
randomized'' concerns an assignment mechanism and its scope. The
assistant should preserve these differences before converting prose into
software constraints.

The same caution applies to exclusions. Absence of an association in one
study can reflect imprecision, cancellation, measurement error or a
different population. Absence of a retrieved passage can reflect search
failure. Neither observation automatically justifies a forbidden causal
path. A hard exclusion removes models from the admitted class; it is
therefore a scientific restriction with consequences for every
downstream identified set.

Let \(\mathcal M\) be the currently admitted set of causal models and
\(K\) a proposed restriction. The restricted class is

\[
 \mathcal M_K=\{m\in\mathcal M:m\text{ satisfies }K\}.
\]

For a fixed observed law and target, its compatible effect set is
contained in the corresponding unrestricted set. Narrowing follows from
set inclusion, not from new empirical accuracy. If the true model is
excluded by \(K\), the narrower result can be more decisive and less
defensible at the same time. If the restriction leaves no compatible
model, the output should identify the contradiction rather than
substitute an empty effect with the meaning ``no effect.''

A useful evidence record therefore includes the exact proposition,
variable instances, population and period, primary source passage,
whether the assertion is observed or assumed, and counterevidence. It
also records the target calculations that depend on the proposition. If
two experts disagree about a common cause, the analyst should be able to
see the effect bounds and feasible actions under each branch. If they
disagree about a remote edge that changes neither the relevant
functional nor the feasible decision, the next review can focus
elsewhere without declaring the remote edge resolved.

One source can support several claims without providing several
independent pieces of evidence. A guideline, a review quoting that
guideline and three generated summaries of the review may all trace to
one original study. The proposed evidence graph should retain that
common origin. It should also preserve differences between a source's
measured population and the bank's target population. Authenticity,
applicability and inferential strength are separate properties of the
same document.

\hypertarget{agreement-between-agents-does-not-measure-causal-certainty}{%
\section{Agreement between agents does not measure causal
certainty}\label{agreement-between-agents-does-not-measure-causal-certainty}}

Consensus can be useful for finding inconsistent explanations, but its
interpretation depends on how the answers were generated. Let
\(V_1,\ldots,V_n\) be binary endorsements of one proposed edge. In a
stipulated exchangeable response model, each has mean \(q\), variance
\(q(1-q)\), and common pairwise correlation \(\rho\). Expanding the
variance of the mean gives

\[
 \operatorname{Var}(\bar V)
 =\frac{q(1-q)}{n}\{1+(n-1)\rho\}.
\]

The variance-equivalent count of independent responses is

\[
 n_{\rm eff}=\frac{n}{1+(n-1)\rho}.
\]

For ten responses and an assumed correlation of .9, this count is
\(100/91\), barely more than one. This is an exact implication of the
response model, not an estimated correlation or a calibration result for
any deployed agent. Shared training, retrieved passages or an initial
proposed answer can create dependence whose size must be studied rather
than guessed. \citep{exec13_diagram_examples}

Even \(\rho=0\) would not solve the more fundamental problem. The
parameter \(q\) describes the chance that the prompted system endorses
an edge. It is not the probability that the edge is true in the target
bank. Independent repetitions can estimate a consistently wrong
endorsement probability very precisely. Assigning different role names
to the same model changes the prompts; it does not establish either
statistical independence or independent access to scientific evidence.

The hallucination-reduction study by Sng and colleagues makes a related
distinction concrete. Its contradiction metric compares model ratings,
its retrieval experiment supplies a mocked true causal graph, and its
debate setup uses opposing positions and an arbiter. The authors
explicitly distinguish reduced contradiction from accurate causal
prediction and leave true-graph validation to further work. That is
useful evidence about a proposed consistency remedy, with a narrower
claim than the paper title might suggest. \citep{exec13_hallucination}

The proposed product can use debate to expose an omitted assumption, ask
for a counterexample or compare the scope of two sources. It should
retain the actual reasons and original passages. A final arbiter's vote
should remain a workflow decision or an explicitly uncalibrated
preference. If the intended downstream operation needs probabilities
over models, those probabilities require a declared prior/likelihood
construction or a separately justified calibration target. A voting
fraction cannot silently supply them.

\hypertarget{reuse-decisions-for-three-available-package-paths}{%
\section{Reuse decisions for three available package
paths}\label{reuse-decisions-for-three-available-package-paths}}

The inspected pywhy-llm revision provides useful hypothesis-elicitation
components. Its identification suggester delegates a back-door
suggestion to the model's confounder suggester, while the inspected
front-door method in that identification-suggester class is a TODO with a
pass body and the longer estimand-construction block is commented out. The
separate simple suggester can generate front-door names, as the preceding
formulation chapter explains; textual suggestions still require formal
identification review. Its MIT license and
declared Python/LLM dependencies make reuse of selected interaction
ideas plausible. They do not make a returned confounder list an
identified adjustment set. The proposed integration should accept
suggestions as candidate evidence records and send the resulting
reviewed graph to a formal identification operation. It should not
present the front-door placeholder as an available scientific
capability. \citep{v12_la_package_pywhy_llm}

The pinned causal\_GPT source is valuable as an experimental comparison
for response-based discovery. It aggregates model answers through
voting, weighted voting and a proportion test, then maps the selected
decision to zero or one in the conditional-independence interface. Those
endpoints are software signals; they are not calibrated p-values for the
bank's joint distribution. The source also reuses YES and NO counts from
the same answer collection, which makes their covariance relevant even
to the response-level test. The analysis in
Chapter~\ref{ch:la03-discovery} develops that separate
issue. The appropriate decision is to retain the method as a labeled
benchmark candidate, with its BSD three-clause obligations and actual
dependencies, while declining direct use as a calibrated sampling-law
oracle. \citep{v12_la_package_causal_gpt}

CausalAgent offers an end-to-end application and MCP comparison. At the
inspected revision, its server advertises PC and OLC tools. It imports
the PC analysis function, while the OLC function called in that module
is neither imported nor defined there; the surrounding exception handler
returns a JSON failure string. A separate output helper takes the text
between the first and last braces and returns an empty dictionary when
decoding fails. These static findings constrain adoption of the
inspected path. They are not a claim that every revision or every
deployment behaves identically. The proposal should reuse the idea of
explicit tool handoffs, require typed invalid/unavailable outcomes, and
defer the OLC path until its concrete implementation is repaired and
actually exercised. \citep{v12_la_package_causalagent}

There are three distinct reuse decisions here. Selected elicitation
components can be integrated behind a proposal-only boundary. An
imperfect scientific method can remain a useful experimental comparator.
A concrete incomplete tool path can be rejected for direct adoption
while retaining other architectural lessons. Porting an entire
application or dependency stack does not resolve any of the scientific
defects. A port should be considered only after a selected path has a
justified purpose and a measured constraint that the port could address.

Actual deployment would still require frozen package and model
identities, license notices, data-handling decisions, runtime budgets,
endpoint conversion checks and independent reference cases. The current
comparison is based on saved source and primary text. No live-model
accuracy, end-to-end latency or runtime conformance is inferred from
that inspection. These limitations define the first prototype trials
without blocking the present architecture and literature argument.

\hypertarget{editing-a-graph-without-losing-its-scientific-history}{%
\section{Editing a graph without losing its scientific
history}\label{editing-a-graph-without-losing-its-scientific-history}}

A graph edit should change one declared object and make its consequences
inspectable. Adding a latent node is different from adding an observed
column. Resolving a circle endpoint is different from moving a node on
the screen. Accepting an expert exclusion is different from an algorithm
finding an orientation under a registered test. The interface should
retain the origin and logical status of each change even when it
displays the resulting graph in one view.

Suppose the analyst begins with a partially directed graph and a
disputed assertion that a risk score is sufficient for adjustment. A
useful edit sequence first records the score as a measured proxy, then
creates a branch with the latent common cause, and finally compares the
target functional under the original and revised assumptions. The proxy
example shows why the changed value can be scientifically meaningful
even when model fit remains good. The analyst should be able to return
to the earlier branch and see the exact premise responsible for the
difference.

Constraint checking also needs a clear purpose. Rejecting cycles
maintains the declared DAG representation, but greedy rejection can
depend on insertion order. Given proposed arrows \(A\to B\), \(B\to C\)
and \(C\to A\), accepting the first two rejects the third; another order
rejects a different arrow. The algorithm has selected one internally
consistent subset without establishing which arrow was scientifically
wrong. The remedy is to expose the conflicting set, preserve the
evidence for each member and seek a revised clock, representation or
premise. A deterministic insertion order improves reproducibility; it
does not supply a causal reason for the choice.

The downstream dependency record should bind an identified expression to
the graph, query and assumptions that produced it. If an edit changes
the outcome's definition, the target changes. If an edit changes a
common-cause restriction, identification may change. If only a displayed
label changes through a valid bijection, the scientific object may
remain the same. These cases deserve different invalidation rules.
Reusing an old estimate because the edited diagram still has the same
number of edges is not a defensible rule.

Review should preserve disagreement when available evidence cannot
resolve it. A graph collection is not automatically a statistical
confidence set, and a union of its effect candidates is not
automatically exhaustive over every causal model. The product must state
which alternatives were considered, why others were excluded and what
additional uncertainty remains. A well explained partial answer can
still narrow the next experiment or rule out a premature investment.

\hypertarget{which-disagreement-is-worth-resolving}{%
\section{Which disagreement is worth
resolving}\label{which-disagreement-is-worth-resolving}}

The monetary question can organize expert review without pretending to
settle causal truth. Consider an assumed proposal scenario in which two
admitted model branches imply effect values .02 and .08 on repayment
probability. Assign an assumed USD 1,000 value to an additional
successful repayment and a USD 50 offer cost. The corresponding net
values per account are

\[
 v_L=1000(.02)-50=-30,\qquad
 v_H=1000(.08)-50=30.
\]

With no justified probability over the two branches, the action is not
robustly beneficial. Relative to choosing the better action in each
branch, both adopting and declining have maximum regret USD 30 per
account. The calculation locates a consequential disagreement; it does
not choose a probability because a majority of prompted agents favored
one branch. \citep{exec13_diagram_examples}

If a sponsor explicitly adopts a planning probability \(\pi\) for the
high-value branch, the expected net value of adoption is \(60\pi-30\).
The best current expected value is \(\max(0,60\pi-30)\). Perfect
information about the branch would allow adoption only when its value is
positive, yielding expected value \(30\pi\). The expected value of
perfect information is therefore

\[
 \operatorname{EVPI}(\pi)
 =30\pi-\max(0,60\pi-30).
\]

At \(\pi=1/2\), it is USD 15 per account; at \(\pi=4/5\), it is USD 6.
For 10,000 otherwise comparable accounts with additive values and one
shared model uncertainty, those assumed upper bounds are USD 150,000 and
USD 60,000. The higher prior confidence reduces the incremental value of
learning even though it makes immediate adoption more attractive. These
are exact planning calculations, not observed bank returns or elicited
institutional beliefs. \citep{exec13_diagram_examples}

A real experiment provides imperfect information, so its expected value
of sample information cannot exceed the perfect-information value under
the same correct decision model. Delay, implementation, review and
measurement costs reduce its net value further. An experiment costing
more than this upper bound is not justified by this particular decision
alone. A lower cost is necessary but not sufficient: one must still show
that the experiment can distinguish the relevant branches and that its
result can change an available action.

A public benchmark can establish that the prototype preserves branch
identities, computes the right functional in a known mechanism and
avoids the proxy/collider mistakes. It generally cannot reveal which
causal mechanism operates in a specific bank. The proposal should
therefore sequence the investment: first test the accessible
computational and reasoning obligations; then ask whether the remaining
target-specific uncertainty is valuable enough to investigate. Expensive
authentic evidence is a conditional next investment, not a prerequisite
for writing a defensible proposal.

\hypertarget{a-benchmark-for-diagram-assistance-with-explicit-failure-credit}{%
\section{A benchmark for diagram assistance, with explicit failure
credit}\label{a-benchmark-for-diagram-assistance-with-explicit-failure-credit}}

Evaluation should begin with a declared task object. One case might ask
for a diagram consistent with a timing description; another for possible
omitted variables; a third for an intervention answer over unresolved
graph alternatives. They need different references. An author-stipulated
graph can provide controlled truth for a generated case. A public
real-world dataset may provide an observation law and metadata without a
unique causal graph. Scoring both against one arbitrarily selected edge
list would confuse missing information with model error.

The proxy and collider cases provide exact formal baselines. The
reference is the stated structural law and its separately enumerated
intervention distribution. The assistant sees the permitted variable
descriptions and observational information, with the hidden generator
truth withheld whenever the task concerns discovery. A correct output
should distinguish what follows from those observations from what
depends on an additional assumption. A refusal to declare
unconfoundedness is useful when the evidence does not supply it; a
refusal to do any arithmetic in a fully specified exact model is a
different failure.

The comparison should include a deterministic dictionary-and-clock
validator, an unassisted formal method in its admitted domain, a
single-model elicitation route, a retrieval-assisted route and a
reviewed human workflow where feasible. Public or generated cases permit
these designs without requiring a bank deployment. A live-model
comparison additionally needs actual model execution, recorded settings
and measured resource use. An offline schema or literature synthesis
supplies neither live accuracy nor superiority.

Perturbations should test meanings that ought to remain invariant.
Permuting variable order should not change the meaning of conditional
independence. Renaming variables through a bijection should preserve a
formal answer when their descriptions remain equivalent. Conversely,
changing a variable from pre-offer state to post-offer response should
change the admissible reasoning. Removing a cited premise should weaken
the claims that depended on it. A benchmark that rewards invariance to
every wording change would miss this distinction between irrelevant
phrasing and substantively changed information.

Scores should retain unsupported, malformed and unavailable requests in
the denominator. They should distinguish valid abstention from an
uninformative missing answer, graph syntax from target identification,
and a justified ambiguous effect set from an omitted calculation. The
monetary example supplies an additional loss measure: did the assistant
recommend an action that the admitted models do not justify, or identify
a potentially valuable next information request? Edge accuracy remains
useful, but it is insufficient when a small number of consequential
edges determines the decision.

Finally, reference information and retrieval material must be separated
deliberately. Giving a model the true graph can test whether it uses
that information correctly, as in a controlled consistency task. It
cannot also measure discovery of that graph from evidence that did not
contain it. The proposed evaluation should report which service was
actually tested, preserve unfavorable cases, and measure the full cost
of retrieving, checking and acting on an answer. That produces a
reviewable reason to build or narrow a diagram-assistance prototype
while leaving target-bank scientific truth to the later study that could
establish it. \label{exec13:diagram-development-end}


% SOURCE docs/release/2026-10-29/source/docs/writing-repair/2026-10-29/src/dc405a80e39693a7.tex CHAR 0
\hypertarget{an-operational-ontology-for-four-banking-journeys}{%
\section{An operational ontology for four banking
journeys}\label{an-operational-ontology-for-four-banking-journeys}}

An ontology is useful when it makes different interpretations explicit
before they enter an estimate. The proposed causal ontology should
distinguish an entity from an observation about that entity, a concept
from a measured variable, a variable definition from a particular
observation, and an intervention from the administrative event used to
record it. These distinctions are not cosmetic. They determine which
quantities can be assigned, which can be conditioned on, which share
measurement errors and which belong to the same population.

For the funding journey, consider ``interest rate.'' A contractual
annual rate, a monthly interest payment, a competitor's advertised rate
and a fitted marginal funding cost are different variables. A unit
conversion can relate two encodings of the same rate, but it cannot turn
a payment into an offer without account balance, accrual and timing
information. Nor does it turn an advertised offer into a received offer.
The ontology should identify the product, currency, contractual basis,
decision date and exposure window. An agent can use a glossary to
suggest the mapping; the analyst must still see which distinctions the
mapping retains.

The associated entity hierarchy also matters. An account belongs to a
legal customer, a customer may belong to a household, and several
records may refer to the same account. Those relations do not themselves
say that one node causally influences another. They do affect
aggregation and dependence. An offer randomized by household cannot be
analyzed as though each account were assigned independently merely
because the extract has one row per account. If a latent household
membership is replaced by a record-linking prediction, that substitution
introduces the same kind of measurement question developed in the proxy
example.

For sequential servicing, ``treatment'' may mean eligibility,
assignment, receipt, completion or adherence. The stages can occur on
different dates and have different causes. A rule that assigns
assistance after an arrears flag targets a dynamic regime; comparing
customers who eventually complete assistance targets a selected group
unless further assumptions justify a different interpretation. The
ontology should retain the action that the bank can choose and the later
events that it can only influence. Otherwise an assistant may propose an
infeasible intervention on a customer's future response while describing
it as an ordinary policy choice.

For discovery, a variable needs a measurement protocol as well as a
label. A repeated risk score can change because the customer changed,
because the scoring model was updated or because new source information
became available. Treating all vintages as the same mechanism can make
an apparent causal change coincide with a measurement change. A diagram
can represent the score-generating mechanism or restrict the analysis to
a stable declared regime. The correct choice depends on the question; a
vocabulary match alone cannot choose it.

For trustworthy observations, separate an event from its record, a
proposed match from a true identity, and label correctness from the
value being labeled. A classifier's probability that a record pair
matches is an algorithm output with a calibration question. It is not an
observed identity. A hand-reviewed match can supply additional
information, but its reviewer, selection process and uncertainty remain
relevant. The ontology can preserve these distinctions in a public
synthetic example long before the proposal needs an expensive bank
labeling exercise.

Across all four journeys, the proposed variable definition therefore
contains a conceptual role, unit of analysis, value domain, measurement
mechanism, clocks, lineage and relation to a feasible intervention. The
system should support unresolved fields explicitly. A missing clock
should prompt a targeted evidence request; an unresolved proxy relation
should produce alternative measurement models. Automatically completing
every field with plausible text would hide the very uncertainty that a
causal proposal must expose.

\hypertarget{semantic-relations-need-typed-boundaries}{%
\section{Semantic relations need typed
boundaries}\label{semantic-relations-need-typed-boundaries}}

Common ontology relations carry different logical meanings. ``A savings
account is an account'' describes a taxonomy. ``A household owns an
account'' describes a relation between entities. ``A risk score is
computed from balance history'' describes a derivation. ``An offer
changes a future balance'' proposes a causal mechanism. A general
relation extractor may return all four as triples, but a causal engine
cannot assign them all the same edge semantics.

A taxonomy helps establish which definitions apply to an instance; it
does not imply that intervening on an instance's category has a coherent
economic meaning. Ownership determines which exposures should be
combined and may enter a separate causal model, but the ownership triple
alone does not identify an intervention effect. A derived score carries
deterministic or stochastic dependence on its inputs; treating the score
and every input as unrelated primitive variables can contradict the
declared derivation. Each relation needs a type before a diagram
operation acts on it.

Duplicate encodings provide a simple formal example. Suppose a decimal
rate \(R\) and its basis-point encoding \(B\) satisfy \(B=10{,}000R\)
exactly. A vocabulary matcher may list them as two numerical variables.
Their observational distribution lies on a deterministic line. An
imagined intervention setting \(R=r\) and \(B=b\) independently is
incompatible with the encoding relation unless \(b=10{,}000r\), or
unless the target is explicitly a system fault that breaks that
relation. Ordinary policy analysis should treat these as two coordinates
of one declared quantity, not as two independently selectable business
levers.

Alias resolution must be reversible. Two fields with the same display
name can describe different vintages; two differently named fields can
be identical encodings. The system should retain source identifiers,
proposed canonical identifiers and the transformation between them. A
bijective unit conversion can be checked mechanically. A many-to-one
coarsening, such as mapping several delinquency states into ``at risk,''
discards information and can change the adjustment or decision problem.
It therefore needs a scientific interpretation as well as a data
transformation.

The preferred local baseline is a deterministic dictionary and lineage
validator. It can detect an undeclared conversion, missing unit,
duplicate identifier or a cycle in a purported derivation. A language
assistant adds value if it finds plausible semantic conflicts that the
dictionary did not encode and supplies reviewable evidence for resolving
them. The two should be compared on missed consequential distinctions
and unnecessary review work. Merely producing more complete metadata is
not enough when some of the added metadata is wrong.

The Selby study provides relevant adverse evidence for automatic
completion. Its inspected experiments elicited quantitative priors and
imputed tabular values from language models. Repeated prompts often
produced consistent responses, yet some priors conflicted strongly with
observations and conventional nearest-neighbor or random-forest
imputation often performed better. The authors identify ambiguous units,
incomplete descriptions and recognizable public datasets as limitations.
This supports careful metadata and benchmark design. It does not
establish a theorem about causal identification or license silent
imputation of a bank's missing assignment mechanism.
\citep{v12_la_2402_07770}

\hypertarget{a-treatment-label-can-hide-several-interventions}{%
\section{A treatment label can hide several
interventions}\label{a-treatment-label-can-hide-several-interventions}}

Consider an assumed assistance program with two versions. Without
assistance, successful repayment has probability .4. The low-intensity
version has success probability .42 and costs USD 10; the high-intensity
version has success probability .6 and costs USD 100. Suppose assignment
and version selection are randomized in this constructed example, and
each additional successful repayment is assigned a planning value of USD
1,000. The version probabilities define the program being evaluated.

If \(w\) is the probability of the high-intensity version among treated
accounts, the effect of that mixed program is

\[
 \tau(w)=(1-w)(.42)+w(.6)-.4=.02+.18w.
\]

Its expected cost is \(c(w)=10+90w\), so the assumed net value is

\[
 v(w)=1000\tau(w)-c(w)=10+90w.
\]

At \(w=1/4\), the effect is .065 and the net value is USD 32.50. At
\(w=3/4\), the effect is .155 and the net value is USD 77.50. Every
calculation is exactly reproduced from the stipulated rational
probabilities. No observed banking profitability is asserted.
\citep{exec13_ontology_examples}

Both datasets could contain the same binary column named ``assistance.''
Nevertheless, the programs have different meanings and different costs.
An assistant should ask whether the deployment preserves the version mix
and the mechanisms under which those version-specific outcomes were
obtained. Reporting a precisely estimated effect for the earlier mix as
the effect of the later mix would answer the wrong question even if the
graph's two main node labels were unchanged.

The construction is deliberately simple. If version choice depends on
baseline need, a dynamic or covariate-dependent mixture \(w(x)\) is
relevant. If customers select their own version after assignment,
assignment and receipt become different interventions. If staff
implement nominally identical assistance differently, the declared
version may still omit consequential variation. A useful ontology does
not promise to remove all such ambiguity; it identifies which variation
is included in the intervention definition and which is left as an
unresolved transport assumption.

This also explains why a prototype should retain version information
before it is needed by an estimator. A source package might accept one
binary treatment and correctly estimate the target supplied to it. The
proposal's common layer is responsible for making that supplied target
faithful to the actual decision. An LLM can identify the hidden-version
possibility and help locate policy documentation. The effect arithmetic
then belongs to a separately specified calculation, with the assumed
assignment and valuation law visible to the reader.

\hypertarget{a-restricted-ontolearner-integration}{%
\section{A restricted OntoLearner
integration}\label{a-restricted-ontolearner-integration}}

OntoLearner supplies an available component for the semantic stage. Its
paper describes ontology loading, learning tasks and learner models,
including typing, taxonomy and relation discovery. It also reports
domain-dependent error patterns, retrieval limitations and calibration
differences. The inspected requirements were defined by the developers,
and counts of ontologies or requirements differ across passages. A
proposal can use the architecture and reported limitations without
turning those counts into a certified inventory or treating every
extracted relation as causal. \citep{v12_la_2607_01977}

The inspected repository revision is
074a41ba23a85f5b72d4e28da042eab3d50bd90a, dated 3 September 2026. Its
loader chooses a local source when a path is supplied and otherwise
calls a Hugging Face download. The inspected download call supplies a
repository, file and dataset type, but no revision. Further extraction
calls download typing, taxonomy and non-taxonomic-relation files without
a revision argument as well. Pinning this Python source therefore does
not by itself pin the ontology bytes that an automatic download will
return. \citep{v12_la_package_ontolearner}

There is another relevant boundary in the imported ontology closure. The
base class's \texttt{contains\_imports} hook returns false. When a
subclass enables that hook, the loader follows resolved imports, but it
catches import-loading failures, logs them and continues; an unresolved
import also produces a warning without failing the load. A nonempty
top-level graph therefore does not establish that all intended imported
definitions were loaded. If an unavailable import contains the
definition distinguishing a legal customer from a household, the
resulting partial ontology could still look operationally successful
while supporting an incomplete semantic map. This is a source-inspected
failure possibility, not a report of an executed OntoLearner trial.
\citep{exec13_ontolearner_inspection}

The proposed integration should admit a fixed local ontology bundle with
explicit revision, byte identities and an import manifest. A wrapper
should report complete, intentionally partial or failed loading,
preserving which definitions are absent. It should carry the original
relation types into a semantic proposal record and require a separate
reviewed mapping before creating causal nodes or edges. The initial
useful output is a candidate variable dictionary with evidence and
conflicts, not an automatically accepted causal graph.

The package's MIT license permits reuse subject to its notice and
disclaimer conditions. Its inspected project file declares Python at
least 3.10 and below 3.14, together with RDF, graph, dataframe,
validation and model-related dependencies; several version
specifications are broad. Reusing a small local loading/extraction path
should still be assessed against a locked actual environment and its
transitive dependencies. A source inspection provides no build, memory
or compatibility measurement. The proposed design is to evaluate a restricted, reproducible semantic path
and defer automatic model-backed ontology expansion until its task, data
rights, resource use and evidence handling have been assessed. The preceding
ontology chapter also shows why document identity must be checked separately
from lexical similarity before its extraction metric can support a causal
observation contract.
\citep{v12_la_package_ontolearner}

Porting the whole library would preserve neither ontology identity nor
causal meaning automatically. A port could be justified later by a
measured platform constraint, but it would require equivalent parsing,
import handling and relation semantics. Direct adoption of a mutable
remote default would make a reproduction depend on an unspecified
ontology version. Rejecting the entire library would discard useful
semantic infrastructure. The restricted path retains its useful
operations while making their scientific handoff explicit.

\hypertarget{evidence-extraction-as-a-separate-experiment}{%
\section{Evidence extraction as a separate
experiment}\label{evidence-extraction-as-a-separate-experiment}}

A claim about an edge has at least three evidential questions. Was the
relevant source obtained and correctly parsed? Does an inspected passage
actually support the narrow proposition attributed to it? Is that
proposition applicable and adequate for the intended banking model? A
citation attached to a graph is useful only if the workflow can answer
these questions separately. Text extraction accuracy cannot establish
causal validity, and a valid causal result from another population does
not automatically establish transport.

The ARGUS study provides a concrete adjacent comparison. It audits
reported evidence for identifying assumptions in
difference-in-differences studies and explicitly does not adjudicate
causal truth. Its flaw-injection experiments supply controlled local
references, while its small human-label pilot diagnoses failures in
real-paper auditing. The gated pipeline's omissions include cases in
which retrieval found no relevant passage and the system abstained. In
the pilot, human reviewers often found evidence for cells where the
system had abstained; the paper's subsequent calibration changes were
evaluated on the same papers. These details matter more to our design
than a single aggregate agreement number. \citep{v12_la_2609_30867}

For a diagram assistant, a suitable first experiment can be public and
inexpensive. Assemble passages whose scope is explicit, paired with
propositions that they support, contradict or leave unresolved. Include
genuine parse failures, omitted context, differing populations and
statements that concern association rather than intervention. An
independent reference should record the exact passage and the reason for
its label. Generated alterations can expose specific failure mechanisms,
but their artificial nature should remain visible and should not be
counted as representative bank documents.

The baseline should include direct passage lookup and a deterministic
field/rule check before introducing a model. A model-based system should
report the retrieved material and the proposition it inferred from it.
If the reference passage was never retrieved, the retrieval component
failed. If it was retrieved but misread, the interpretation failed. If
it was read accurately but promoted beyond its population or
assumptions, the scientific mapping failed. Different remedies follow:
improve acquisition/search, improve task interpretation, or tighten the
rules for accepting a causal premise.

An observed negative result also needs precise wording. ``No relevant
evidence was found in this bounded search'' records an acquisition
outcome. ``The inspected document does not state how assignment
occurred'' records a document-level absence under the inspected scope.
``Assignment was confounded'' is a claim about the underlying process.
The first two do not prove the third. Conversely, a polished statement
that assignment was ``as if random'' does not establish the required
condition merely because the text contains the expected vocabulary.

This experiment can support a decision about building a literature and
documentation assistant before paying for a bank audit. It cannot show
that the bank's declared assignment process is true. The eventual
external input is a process-specific record or experiment, sought only
if the prototype and the investment calculation establish its value. The
staged evidence argument therefore has a concrete boundary: public
documents test extraction and reasoning obligations; target-specific
evidence tests propositions about the target institution.

\hypertarget{reviewing-conflicts-without-counting-votes-as-evidence}{%
\section{Reviewing conflicts without counting votes as
evidence}\label{reviewing-conflicts-without-counting-votes-as-evidence}}

Consider three separately proposed restrictions: an offer changes
engagement, engagement changes a later score, and the score determines
the earlier offer. They cannot all be direct temporal assertions about
the same variable instances in an acyclic model. A syntax checker can
identify the contradiction. It cannot determine whether the score was
actually measured earlier, whether ``changes'' meant association,
whether a date was omitted, or whether a cyclic equilibrium model was
intended. The review should return these alternatives to the analyst
with their evidence.

One useful formal object is the set of admitted models after each
restriction. If \(K_1,K_2,K_3\) are proposed predicates, the combined
class is

\[
 \{m\in\mathcal M:K_1(m)\land K_2(m)\land K_3(m)\}.
\]

If it is empty, the system can search for a smallest subset of the
predicates that is already inconsistent within the declared model class.
Reporting that subset localizes the conflict. It does not rank the
predicates by truth or show that the smallest deletion is scientifically
appropriate. Several minimal conflicting subsets may exist; suppressing
all but the first can leave the analyst unaware of another
contradiction.

Priority should follow the evidence and the decision. A verified
timestamp can refute a purported same-instance temporal order. An
expert's general mechanism description may instead need to be
instantiated at a different date. A model suggestion with no supporting
source may remain a hypothesis. If withdrawing one disputed restriction
changes a robustly positive value into a set containing a substantial
loss, that restriction merits attention before a remote relation whose
alternatives leave the decision unchanged. This is a proposed review
policy, not a claim that the system can mechanically rank every kind of
scientific evidence.

Conflicts between sources also need not be contradictions. Two studies
can investigate different treatment versions, observation regimes or
populations. The ontology and intervention-version records make those
differences explicit. A source that supports an association at branch
level may have little to say about an account-level intervention. A
diagram can preserve both claims at their proper levels rather than
force one to win a debate about an underspecified phrase.

The proposal's deliverable is thus a reviewable sequence from language
to typed variables, explicit premises, compatible models, target
calculations and decision consequences. Each transition has examples
that can pass, fail or remain unresolved. That sequence gives an agent
useful work throughout the causal analysis while preserving the
mathematical and evidential limits of each contribution.
\label{exec13:ontology-development-end}


% SOURCE docs/release/2026-10-29/source/docs/writing-repair/2026-10-29/src/266aba6506dc1143.tex CHAR 0
\hypertarget{why-an-accurate-clock-need-not-identify-the-answer}{%
\section{Why an accurate clock need not identify the
answer}\label{why-an-accurate-clock-need-not-identify-the-answer}}

It is useful to follow one restriction through an entire finite model
comparison. Return to an offer \(A\) and a later repayment outcome
\(Y\), both binary. Consider three explicitly listed models, using
mutually independent primitive variables in each model. In model \(F\),
\(A\) is fair Bernoulli and \(Y=A\mathbin{\oplus}E\), where \(E\) is
Bernoulli with probability \(1/10\) and \(\oplus\) denotes exclusive OR.
In model \(R\), \(Y\) is fair Bernoulli and \(A=Y\mathbin{\oplus}E\). In
model \(C\), a pre-offer latent variable \(U\) is fair Bernoulli,
\(A=U\mathbin{\oplus}E\) and \(Y=U\). Each error variable has the same
probability \(1/10\). These are alternative constructions, not
simultaneous equations for one population.

Every model produces the observed table

\[
 P(A=0,Y=0)=P(A=1,Y=1)=\frac9{20},\qquad
 P(A=0,Y=1)=P(A=1,Y=0)=\frac1{20}.
\]

All four observed cells are positive. In model \(C\), even conditioning
on the latent variable leaves both offer assignments possible:
\(P(A=1\mid U=0)=1/10\) and \(P(A=1\mid U=1)=9/10\). Thus the comparison
does not rely on an empty observed treatment stratum or deterministic
assignment given the confounder. Deterministic repayment given \(U\) is
part of this constructed example, and is not a banking empirical claim.

Replacing the action equation gives mean outcomes \(1/10\) and \(9/10\)
in model \(F\), hence effect \(4/5\). In \(R\), replacing the action
equation leaves the fair outcome unchanged; in \(C\), it leaves the fair
latent variable and \(Y=U\) unchanged. Both have effect zero. A separate
exact enumeration verifies the common observed table, positive
assignment probabilities and intervention contrasts.
\citep{exec13_diagram_restrictions}

Now introduce a verified restriction: this particular offer was chosen
before this particular repayment outcome was generated, with no
backward-time causal mechanism in the admitted class. Model \(R\)'s
direct arrow from this outcome to the earlier offer is incompatible with
that restriction. Models \(F\) and \(C\) both respect it. The candidate
graph list shrinks from three models to two, but the candidate effect
set stays

\[
 \{0,4/5\}.
\]

Correctly resolving an orientation has not reduced the decision-relevant
ambiguity in this example. An edge score that rewards removal of \(R\)
can therefore improve while an effect-set or decision measure remains
unchanged. Both statements are useful; they measure different progress.
The set above belongs only to the three-model comparison. It is not
asserted to exhaust every effect compatible with the observed law under
an unrestricted hidden-variable model.

There are several possible next restrictions, with different evidential
costs. If the offer equation is known to use independent random
assignment, the stipulated mechanism of \(C\) is excluded too; among
these candidates only \(F\) remains. If \(U\) becomes correctly
observed, adjustment can identify the zero effect within \(C\)'s
declared mechanism. If a noisy measurement of \(U\) becomes available,
the earlier proxy example shows why adjustment needs further
justification. These routes should not be compressed into one label such
as ``expert-validated diagram.'' They add distinct information and
impose distinct assumptions.

An offer timestamp alone supplies none of the exogeneity information. A
policy document describing randomization might support it, but the
analyst must establish that the document applies to the relevant
assignment, population and period. A random seed recorded after an
officer overrode assignment does not establish the same mechanism. A
public controlled experiment can test whether software applies the
restriction correctly; it cannot verify an unobserved institution's
compliance with its policy. That is the later, targeted evidence
question.

The assistant should show a consequence table after each accepted
premise: which models were removed, which target values were removed,
and which claims now depend on that premise. A useful diagnostic
temporarily withdraws the premise and recomputes those consequences.
Here withdrawing the time order restores \(R\) without changing the two
effect values, while withdrawing the additional exogeneity restriction
restores a zero-effect alternative. This sensitivity report makes the
pivotal premise explicit and can guide the next experiment without
assigning invented probabilities to the models.

\hypertarget{minimal-conflicts-locate-a-question-they-do-not-settle-it}{%
\section{Minimal conflicts locate a question; they do not settle
it}\label{minimal-conflicts-locate-a-question-they-do-not-settle-it}}

The earlier offer--engagement--score example can also be made exact. Let
three restrictions require the direct arrows \(A\to M\), \(M\to S\) and
\(S\to A\) between the same three variable instances in a DAG. All three
together create a directed cycle. Removing any one leaves an acyclic
directed path. The triple is therefore an inclusion-minimal inconsistent
set of these restrictions within the DAG class. A small exhaustive
closure calculation verifies this combinatorial statement.
\citep{exec13_diagram_restrictions}

Here ``inclusion-minimal'' means that deleting any member removes this
conflict. In a larger restriction system it need not mean the smallest
possible cardinality; different minimal conflicting subsets can coexist.
Nor does minimality select a scientifically preferred repair. Deleting
the arrow with the lowest language-model confidence can yield a
syntactically valid graph while deleting the only premise supported by a
reliable process record.

One legitimate repair might refine the score into \(S_t\) and
\(S_{t+1}\). The earlier score affects \(A_t\), which changes engagement
and a later score; that time-unrolled construction can respect all the
underlying mechanism descriptions. Another repair might find that the
phrase ``engagement changes score'' described a deterministic
calculation rather than the primitive intervention contemplated by the
analyst. A third might withdraw an unsupported causal assertion. These
are alternatives to investigate, not guaranteed equivalent versions of
one graph.

The scientific review should therefore report the conflicting predicates
together with their variable identities, clocks and evidence. It should
distinguish a change to the model class, a refinement of a variable, and
a rejection of a premise. Keeping that distinction makes both favorable
and adverse prototype results interpretable: a detector can correctly
find the contradiction even if it cannot resolve its scientific cause.
Its useful contribution is a precise next question whose answer would
change the admissible analysis.

\subsection{The diagram is a set of consequential commitments}

The complete comparison makes graph review a scientific operation. A variable
can be a noisy proxy, a later response or a different version of the intended
action. An accepted edge or exclusion changes the model class; only its effect
on the target and available decisions shows whether that change resolves the
important ambiguity. The retained examples distinguish exact arithmetic under
a stipulated law from evidence that the law applies to a bank.

The useful assistant therefore preserves alternative definitions and graphs,
explains which premises drive each result, and directs review toward
consequential disagreements. Formal methods, deterministic semantic checks and
human evidence review remain explicit alternatives or components. The next
chapter studies language assistance inside discovery procedures, where the
same distinctions must survive statistical tests, graph search and uncertainty
reporting.
\label{exec13:restriction-development-end}

